\documentclass[11pt]{article}

\usepackage[a4paper,margin=1in]{geometry}

\usepackage{booktabs}

\usepackage{longtable}

\usepackage{array}

\usepackage{float}

\usepackage{graphicx}

\usepackage{caption}

\usepackage{adjustbox}

\usepackage{hyperref}

\usepackage{xcolor}

\title{OV$^{2}$SLAM, ORB-SLAM3, and RTAB-Map EuRoC Comparison Report}

\author{Amar Aly}

\date{13-04-2026}

\begin{document}

\maketitle
% Source - https://stackoverflow.com/a/547072
% Posted by godbyk
% Retrieved 2026-04-13, License - CC BY-SA 2.5

\hypersetup{%
    pdfborder = {0 0 0}
}

\tableofcontents

\section{Introduction}

In this report, OV$^{2}$SLAM is benchmarked against ORB-SLAM3 on the EuRoC MAV Sequences. The purpose of this benchmarking is to provide insights on the bottlenecks and accuracy drops because of this setup shift. The benchmarking was done on 10 of the EuRoC MAV (Micro-Aerial Vehicle) sequences, with MH being Machine Hall and V being Vicon Room. Do note that the evaluation setup was done in a way such that it replicates the research papers results.

\textbf{DISCLAIMER}: The experiments in this report are not fully reproducible due to the frequent failures and varying results experienced when running SLAM algorithm. A hypothesis for this would be the different CPU runtime conditions, thread scheduling, missing frames, and non-seeded random actions. Therefore all numbers reported are the \textbf{BEST}-case, because the worst case for these algorithm is an algorithm having 30cm RMSE error and effectively wrong trajectories.
Any sequence red indicates that it failed on ORB-SLAM3 and its number is not representative of the actual performance.

\section{Hardware Details}
All Experiments were ran on the Raspberry PI 5 with 16GB Memory. All Sequences were ran several times and the best RMSE for a sequence is defined as the RMSE that is reasonably close to the research paper, with an important note that many trials got very high RMSE at some times but with rerunning them, the result improved.

\section{Algorithm Implementation Details}

The OV$^{2}$SLAM code was taken from its original source GitHub implementation and benchmarked directly in its ROS2-based form, with only limited changes required to support the experimental setup used in this work. In other words, OV$^{2}$SLAM already provided a ROS-ready code, so it was evaluated more directly than ORB-SLAM3.

However, ORB-SLAM3 was messier as the original source code does \textbf{NOT} provide a ROS implementation, and rather reliance on community-maintained github repositories was necessary. Most repositories are ROS wrappers over the original ORB-SLAM3 code, despite this, the codebase iwht the following URL \url{https://github.com/Mechazo11/ros2_orb_slam3/tree/jazzy} was used as the base ROS2 Jazzy ORB-SLAM3. However, this codebase was only a starting skeleton (according to the github) and required extensive modification to make it suitable for benchmarking in the target environment.

Modifications to ORB-SLAM3 ROS2 were needed in order to allow for smooth benchmarking.
The modifications one where:
\begin{itemize}
    \item fixing hard-coded paths, correcting YAML and package path handling, resolving build and dependency issues related to Pangolin
    \item  timestamp handling had to be corrected from nanoseconds to seconds
    \item fixing the code tendency to use streamed images rather than a bag file, by subscribing to the topic of images
    \item adding Stereo mode, since the current github does \textbf{NOT} include it
    \item inplementing a shutdown ACK-based sequence to allow for automated shutdown instead of manual CTRL+C
    \item disabling autoAbortLBA when a new keyframe arrives because there was a case where the BA completed 44\% of its iterations, which was one of the major causes of high RMSE
\end{itemize}

Overall, the implementation effort differed significantly between the two systems. OV$^{2}$SLAM was \textbf{relatively} straightforward to benchmark because it was already closer to the required ROS workflow, whereas ORB-SLAM3 required a considerable amount of adaptation, debugging, and instrumentation before fair benchmarking could be performed. Do note that more debugging is needed on ORB-SLAM3 as it is considerably less stable than OV$^{2}$SLAM, with sometimes sequences crashing and failing to run. This has been demonstrated on the OV$^{2}$SLAM paper, where they did not show all sequences on ORB-SLAM2 because it consistently failed. One may ask for the relationship between ORB-SLAM2 and ORB-SLAM3, and the answer is
shown in the table below
\begin{table}[ht]
\centering
\caption{Comparison between ORB-SLAM2 and ORB-SLAM3.}
\label{tab:orbslam2_vs_orbslam3}
\renewcommand{\arraystretch}{1.2}
\setlength{\tabcolsep}{5pt}
\begin{tabular}{p{3.2cm}p{5.2cm}p{5.2cm}}
\toprule
\textbf{Aspect} & \textbf{ORB-SLAM2} & \textbf{ORB-SLAM3} \\
\midrule
Release / generation
& Earlier generation system
& Newer generation system, extending ORB-SLAM2 \\
\midrule
Supported visual sensors
& Monocular, stereo, RGB-D
& Monocular, stereo, RGB-D \\
\midrule
Visual-inertial support
& No native IMU support
& Yes, supports visual-inertial SLAM \\
\midrule
Lens models
& Standard pinhole camera model
& Pinhole and fisheye camera models \\
\midrule
Map handling
& Single-map SLAM with map reuse
& Multi-map SLAM with map creation and map merging \\
\midrule
Loop closing / relocalization
& Yes
& Yes, with improved place recognition and map merging capabilities \\
\midrule
Main novelty
& Unified ORB-based SLAM for monocular, stereo, and RGB-D in real time
& Extends ORB-SLAM2 with visual-inertial SLAM, multi-map operation, and stronger robustness \\
\midrule
Claimed accuracy
& Accurate real-time SLAM on standard CPUs
& Reported to be significantly more accurate, with paper claiming about 2$\times$ to 5$\times$ accuracy improvement over prior approaches in some settings \\
\midrule
Typical use case
& Pure vision SLAM benchmarking and deployment
& More advanced SLAM settings, especially when IMU or multi-session robustness is needed \\
\bottomrule
\end{tabular}
\end{table}
\begin{itemize}
    \item ORB-SLAM2 official repository: \url{https://github.com/raulmur/ORB_SLAM2}
    \item ORB-SLAM3 official repository: \url{https://github.com/UZ-SLAMLab/ORB_SLAM3}
    \item ORB-SLAM2 paper: \url{https://arxiv.org/abs/1610.06475}
    \item ORB-SLAM3 paper: \url{https://arxiv.org/abs/2007.11898}
\end{itemize}
\textbf{HOWEVER}, neither algorithms were developed for edge-devices, and failure on edge devices should be expected.
\section{Runtime Configuration}

\subsection{ORB-SLAM3 Setup Details}


\begin{table}[H]
\centering
\small
\caption{ORB-SLAM3 key algorithm parameters from the EuRoC stereo settings file.}
\label{tab:orb_algorithm_params}
\begin{tabular}{lp{9.5cm}}
\toprule
Parameter & Value \\
\midrule
Camera type & PinHole \\
Image size & 752 $\times$ 480 \\
Camera FPS & 20 \\
Color order & RGB (\texttt{Camera.RGB = 1}) \\
Stereo depth threshold & \texttt{Stereo.ThDepth = 60.0} \\
ORB features per image* & \texttt{ORBextractor.nFeatures = 1200} \\
Scale factor & \texttt{ORBextractor.scaleFactor = 1.2} \\
Pyramid levels* & \texttt{ORBextractor.nLevels = 8} \\
FAST initial threshold & \texttt{ORBextractor.iniThFAST = 20} \\
FAST minimum threshold & \texttt{ORBextractor.minThFAST = 7} \\
Left/right topics & \texttt{/cam0/image\_raw}, \texttt{/cam1/image\_raw} via ROS2 wrapper \\
\bottomrule
\end{tabular}
\end{table}
*ORB features per image and Pyramid levels were changed from sequences \textcolor{red}{MH\_04\_difficult} to V2\_02\_medium because of frequent failures. Reasons speculated to be: Low contrast settings requiring more features and more depth to identify changes in frames. Another speculation is that tracking delay was observed to exceed 50ms (< 20fps) in most sequences, which is faster than what the EuRoC streamer tracking speed of 20fps, which caused loss in keyframes and increased RMSE.

\subsection{OV$^{2}$SLAM Setup Details}

\begin{table}[H]
\centering
\small
\caption{OV$^{2}$SLAM key algorithm parameters for accurate and fast stereo configurations.}
\label{tab:ov2_algorithm_params}
\begin{tabular}{lp{4.6cm}p{4.6cm}}
\toprule
Parameter & Accurate stereo & Fast stereo \\
\midrule
Loop closer & \texttt{buse\_loop\_closer = 1} & \texttt{buse\_loop\_closer = 0} \\
Realtime flag & \texttt{force\_realtime = 1} & \texttt{force\_realtime = 1} \\
FAST detector & \texttt{use\_fast = 0} & \texttt{use\_fast = 1} \\
Shi-Tomasi detector* & \texttt{use\_shi\_tomasi = 0} & \texttt{use\_shi\_tomasi = 0} \\
BRIEF descriptor & \texttt{use\_brief = 1} & \texttt{use\_brief = 1} \\
Single-scale detector & \texttt{use\_singlescale\_detector = 1} & \texttt{use\_singlescale\_detector = 0} \\
Max keypoint spacing & \texttt{nmaxdist = 25} & \texttt{nmaxdist = 50} \\
FAST threshold & \texttt{nfast\_th = 10} & \texttt{nfast\_th = 10} \\
CLAHE & \texttt{use\_clahe = 1} & \texttt{use\_clahe = 0} \\
Track local map & \texttt{bdo\_track\_localmap = 1} & \texttt{bdo\_track\_localmap = 1} \\
P3P enabled & \texttt{dop3p = 0} & \texttt{dop3p = 1} \\
Loop / mapping mode & \texttt{slam\_mode = 1} & \texttt{slam\_mode = 1} \\
KLT enabled & \texttt{do\_klt = 1} & \texttt{do\_klt = 1} \\
Keyframe filtering ratio & \texttt{fkf\_filtering\_ratio = 0.95} & \texttt{fkf\_filtering\_ratio = 0.9} \\
Final full BA & \texttt{do\_full\_ba = 1} & \texttt{do\_full\_ba = 0} \\
Stereo rectification & \texttt{bdo\_stereo\_rect = 0} & \texttt{bdo\_stereo\_rect = 0} \\
Undistortion & \texttt{bdo\_undist = 0} & \texttt{bdo\_undist = 0} \\
Image size & 752 $\times$ 480 & 752 $\times$ 480 \\
Topics & \texttt{/cam0/image\_raw}, \texttt{/cam1/image\_raw} & \texttt{/cam0/image\_raw}, \texttt{/cam1/image\_raw} \\
\bottomrule
\end{tabular}
\end{table}
* OV$^{2}$SLAM claims to use Shi-Tomsai, but code is showing that OV$^{2}$SLAM is using a custom keypoint extractor.
\subsection{RTAB-Map Setup Details}

RTAB-Map was evaluated in stereo frame-to-map odometry mode using the \texttt{rtabmap\_ros} package under ROS2 Jazzy. Camera images were first rectified via two \texttt{image\_proc/rectify} nodes and camera info published by a custom \texttt{euroc\_caminfo\_pub.py} node. The \texttt{stereo\_odometry} node performed visual odometry using GFTT feature detection and FLANN-based matching, while the \texttt{rtabmap} node performed global SLAM with appearance-based loop closure via a bag-of-words index. All sequences were replayed from \texttt{.mcap} ROS2 bags without time manipulation (\texttt{bag\_rate: 1.0}); frame dropping is therefore possible when processing exceeds the 50\,ms / 20\,Hz budget.

\begin{table}[H]
\centering
\small
\caption{RTAB-Map key algorithm parameters used for EuRoC stereo evaluation.}
\label{tab:rtabmap_algorithm_params}
\begin{tabular}{ll}
\toprule
Parameter & Value \\
\midrule
Feature detector & GFTT (\texttt{Kp/DetectorStrategy = 1}) \\
Max features per frame & \texttt{Kp/MaxFeatures = 300} \\
GFTT min distance & \texttt{GFTT/MinDistance = 4} \\
Odometry map size & \texttt{OdomF2M/MaxSize = 1000} \\
Visual features (odometry) & \texttt{Vis/MaxFeatures = 600} \\
Odometry strategy & Frame-to-map (\texttt{OdomF2M}) \\
Process most recent frame & \texttt{always\_process\_most\_recent\_frame = false} \\
RTAB-Map time threshold & \texttt{Rtabmap/TimeThr = 0} (disabled) \\
RTAB-Map memory threshold & \texttt{Rtabmap/MemoryThr = 0} (disabled) \\
Subscribe stereo & \texttt{subscribe\_stereo = true} \\
Subscribe depth & \texttt{subscribe\_depth = false} \\
Camera FPS (bag) & 20 Hz \\
Left/right topics & \texttt{/stereo\_camera/left/image\_rect}, \texttt{.../right/image\_rect} \\
\bottomrule
\end{tabular}
\end{table}

\clearpage

\section{Three-way Per-sequence Overview}

\begin{table}[H]
\centering
\small
\caption{Four-way per-sequence RMSE summary. OV$^2$SLAM and ORB-SLAM3 report best-run RMSE (Sim(3) alignment); RTAB-Map reports mean$\pm$std across all valid runs (RMSE $>$\,0.01\,m, $\geq$\,50 poses).}
\label{tab:three_way_summary_rmse}
\begin{tabular}{lrrrr}
\toprule
Sequence & OV2 accurate (m) & OV2 fast (m) & ORB keyframe (m) & RTAB-Map avg (m) \\
\midrule
MH\_01\_easy & 0.040 & 0.054 & 0.043 & $0.102 \pm 0.063$ \\
MH\_02\_easy & 0.049 & 0.052 & 0.071 & $0.190 \pm 0.143$ \\
MH\_03\_medium & 0.046 & 0.079 & 0.247 & $0.146 \pm 0.038$ \\
\textcolor{red}{MH\_04\_difficult} & 0.069 & 0.152 & 0.039 & $0.224 \pm 0.070$ \\
MH\_05\_difficult & 0.067 & 0.142 & 0.091 & $0.166$ (1/3 valid) \\
V1\_01\_easy & 0.055 & 0.094 & 0.131 & $0.096 \pm 0.029$ \\
\textcolor{red}{V1\_02\_medium} & 0.088 & 0.095 & 0.331 & $0.081$ (1/3 valid) \\
\textcolor{red}{V1\_03\_difficult} & 0.118 & 0.670 & 0.069 & $0.133 \pm 0.046$ \\
\textcolor{red}{V2\_01\_easy} & 0.072 & 0.086 & n/a & $0.161 \pm 0.052$ \\
V2\_02\_medium & 0.054 & 0.204 & 0.228 & $0.244 \pm 0.035$ \\
\bottomrule
\end{tabular}
\end{table}

\begin{table}[H]
\centering
\small
\caption{Four-way per-sequence primary latency summary.}
\label{tab:three_way_summary_latency}
\begin{tabular}{lrrrr}
\toprule
Sequence & OV2 acc FE (ms) & OV2 fast FE (ms) & ORB track (ms) & RTAB odom (ms) \\
\midrule
MH\_01\_easy & 17.083 & 4.895 & 50.106 & 101.7 \\
MH\_02\_easy & 17.088 & 5.025 & 48.757 & 104.4 \\
MH\_03\_medium & 19.522 & 5.088 & 51.713 & 97.2 \\
\textcolor{red}{MH\_04\_difficult} & 15.578 & 4.990 & 38.323 & 104.2 \\
MH\_05\_difficult & 15.815 & 5.036 & 46.074 & 88.3 \\
V1\_01\_easy & 19.733 & 5.114 & 32.325 & 103.6 \\
\textcolor{red}{V1\_02\_medium} & 21.698 & 5.698 & 37.758 & 93.0 \\
\textcolor{red}{V1\_03\_difficult} & 21.225 & 6.105 & 46.427 & 92.8 \\
\textcolor{red}{V2\_01\_easy} & 15.033 & 4.908 & n/a & 100.3 \\
V2\_02\_medium & 20.488 & 5.793 & 36.495 & 90.4 \\
\bottomrule
\end{tabular}
\end{table}

\begin{table}[H]
\centering
\small
\caption{Four-way per-sequence mean process CPU utilization summary.}
\label{tab:three_way_summary_cpu}
\begin{tabular}{lrrrr}
\toprule
Sequence & OV2 acc (\%) & OV2 fast (\%) & ORB (\%) & RTAB-Map (\%) \\
\midrule
MH\_01\_easy & 153.4 & 42.8 & 139.6 & 132.3 \\
MH\_02\_easy & 149.6 & 43.6 & 134.5 & 134.6 \\
MH\_03\_medium & 160.9 & 48.4 & 145.5 & 125.3 \\
\textcolor{red}{MH\_04\_difficult} & 147.0 & 42.7 & 149.6 & 130.0 \\
MH\_05\_difficult & 147.5 & 42.3 & 143.6 & 131.2 \\
V1\_01\_easy & 157.1 & 46.1 & 87.4 & 131.1 \\
\textcolor{red}{V1\_02\_medium} & 164.9 & 46.8 & 112.2 & 120.1 \\
\textcolor{red}{V1\_03\_difficult} & 162.0 & 46.6 & 132.0 & 120.3 \\
\textcolor{red}{V2\_01\_easy} & 141.7 & 42.0 & n/a & 128.2 \\
V2\_02\_medium & 163.9 & 51.9 & 118.8 & 123.0 \\
\bottomrule
\end{tabular}
\end{table}

\section{OV$^{2}$SLAM Accurate vs Fast Tables}

\begin{table}[H]
\centering
\small
\caption{OV$^{2}$SLAM accurate stereo vs fast stereo per-sequence RMSE comparison.}
\label{tab:ov2_accurate_fast_rmse_summary}
\begin{tabular}{lrrr}
\toprule
Sequence & Accurate RMSE (m) & Fast RMSE (m) & RMSE delta (m) fast-accurate \\
\midrule
MH\_01\_easy & 0.040 & 0.054 & 0.014 \\
MH\_02\_easy & 0.049 & 0.052 & 0.003 \\
MH\_03\_medium & 0.046 & 0.079 & 0.032 \\
\textcolor{red}{MH\_04\_difficult} & 0.069 & 0.152 & 0.083 \\
MH\_05\_difficult & 0.067 & 0.142 & 0.075 \\
V1\_01\_easy & 0.055 & 0.094 & 0.039 \\
\textcolor{red}{V1\_02\_medium} & 0.088 & 0.095 & 0.006 \\
\textcolor{red}{V1\_03\_difficult} & 0.118 & 0.670 & 0.552 \\
\textcolor{red}{V2\_01\_easy} & 0.072 & 0.086 & 0.015 \\
V2\_02\_medium & 0.054 & 0.204 & 0.150 \\
\bottomrule
\end{tabular}
\end{table}

\begin{table}[H]
\centering
\small
\caption{OV$^{2}$SLAM accurate stereo vs fast stereo frontend latency comparison.}
\label{tab:ov2_accurate_fast_fe_summary}
\begin{tabular}{lrrr}
\toprule
Sequence & Accurate FE mean (ms) & Fast FE mean (ms) & FE speedup \\
\midrule
MH\_01\_easy & 17.083 & 4.895 & 3.49x \\
MH\_02\_easy & 17.088 & 5.025 & 3.40x \\
MH\_03\_medium & 19.522 & 5.088 & 3.84x \\
\textcolor{red}{MH\_04\_difficult} & 15.578 & 4.990 & 3.12x \\
MH\_05\_difficult & 15.815 & 5.036 & 3.14x \\
V1\_01\_easy & 19.733 & 5.114 & 3.86x \\
\textcolor{red}{V1\_02\_medium} & 21.698 & 5.698 & 3.81x \\
\textcolor{red}{V1\_03\_difficult} & 21.225 & 6.105 & 3.48x \\
\textcolor{red}{V2\_01\_easy} & 15.033 & 4.908 & 3.06x \\
V2\_02\_medium & 20.488 & 5.793 & 3.54x \\
\bottomrule
\end{tabular}
\end{table}

\begin{figure}[H]
\centering
\includegraphics[width=\textwidth]{images/ov2_rmse_by_sequence.png}
\caption{OV$^{2}$SLAM accurate vs fast RMSE by sequence.}
\label{fig:ov2_rmse}
\end{figure}

As seen from the plot above, the Fast mode of OV$^{2}$SLAM is consistently worse RMSE than the Accurate mode, with notably one example being \textcolor{red}{V1\_03\_difficult} having very high RMSE compared to the research paper. It can be observed that this sequence \textcolor{red}{V1\_03\_difficult} has very intertwined paths that increase the RMSE and confuse a SLAM system.


\begin{figure}[H]
\centering
\includegraphics[width=\textwidth]{images/ov2_rmse_delta_by_sequence.png}
\caption{OV$^{2}$SLAM RMSE delta by sequence (fast minus accurate).}
\label{fig:ov2_rmse_delta}
\end{figure}

From this plot, we can observe that the differences between the Fast and Accurate modes on the EuRoC dataset is generally small, with the expection of \textcolor{red}{V1\_03\_difficult} and V2\_02\_medium due to its very difficult trajectory (seen below). This further supports the fact that the Fast mode is better suited for edge devices due to its speed and acceptable accuracy loss.


\begin{figure}[H]
\centering
\includegraphics[width=\textwidth]{images/ov2_frontend_latency_by_sequence.png}
\caption{OV$^{2}$SLAM frontend latency by sequence.}
\label{fig:ov2_fe_latency}
\end{figure}

In this plot, we observe the front-end latency when running OV2SLAM on each sequence. The lowest latency always comes from the Fast configuration and it has the least standard deviation compared to the Accurate mode. The least possible total front-end latency is around 5ms which corresponds to 200 Hz in Fast mode, while the Accurate mode has a least possible total front-end latency of around 15ms, which is 66 Hz.

\begin{figure}[H]
\centering
\includegraphics[width=\textwidth]{images/ov2_process_cpu_by_sequence.png}
\caption{OV$^{2}$SLAM mean process CPU utilization by sequence.}
\label{fig:ov2_cpu}
\end{figure}
From this plot, we observe that Accurate mode consistently uses significantly more CPU than Fast mode across all sequences. Accurate mode can reach up to 250\%+ CPU utilization and even higher due to the sudden utilization spikes by the Loop Closure threads. In contrast, Fast mode has much lower CPU utilization and can typically run in < 2 cores, though 2 cores would be ideal. This shows that while Accurate mode maximizes estimation accuracy, it comes at a significant cost for edge devices, compared to Fast mode which can have reasonably closer accuracy and much lower latency (3x times reduction in front-end delay)

\begin{figure}[H]
\centering
\includegraphics[width=\textwidth]{images/ov2_role_means.png}
\caption{OV$^{2}$SLAM logical thread-role utilization: accurate vs fast.}
\label{fig:ov2_roles}
\end{figure}
From this plot, we observe that CPU usage in Accurate mode is dominated by the estimator thread, which is responsible for performing local bundle adjustment. The second most computationally intensive thread is the loop closure thread, which adds a substantial amount of CPU consumption on top of the existing processing threads. In contrast, Fast mode shows significantly lower CPU usage across all threads, reflecting its lightweight design. This reduction in computational load supports the idea that Fast mode is more suitable for edge devices while still maintaining high processing speed.


\clearpage

\section{Paper Comparisons}

\begin{table}[H]
\centering
\small
\caption{OV$^{2}$SLAM accurate stereo vs OV$^{2}$SLAM paper stereo rows with loop closure enabled.}
\label{tab:ov2_accurate_vs_paper}
\begin{tabular}{lrrr}
\toprule
Sequence & Accurate stereo RMSE (m) & Paper OV$^{2}$SLAM w. LC (m) & Delta (m) \\
\midrule
MH\_01\_easy & 0.040 & 0.040 & 0.000 \\
MH\_02\_easy & 0.049 & 0.040 & 0.009 \\
MH\_03\_medium & 0.046 & 0.040 & 0.006 \\
\textcolor{red}{MH\_04\_difficult} & 0.069 & 0.060 & 0.009 \\
MH\_05\_difficult & 0.067 & 0.070 & -0.003 \\
V1\_01\_easy & 0.055 & 0.090 & -0.035 \\
\textcolor{red}{V1\_02\_medium} & 0.088 & 0.070 & 0.018 \\
\textcolor{red}{V1\_03\_difficult} & 0.118 & 0.090 & 0.028 \\
\textcolor{red}{V2\_01\_easy} & 0.072 & 0.070 & 0.002 \\
V2\_02\_medium & 0.054 & 0.060 & -0.006 \\
\bottomrule
\end{tabular}
\end{table}

\begin{table}[H]
\centering
\small
\caption{OV$^{2}$SLAM fast stereo vs OV$^{2}$SLAM paper stereo rows without loop closure.}
\label{tab:ov2_fast_vs_paper}
\begin{tabular}{lrrr}
\toprule
Sequence & Fast stereo RMSE (m) & Paper OV$^{2}$SLAM no LC (m) & Delta (m) \\
\midrule
MH\_01\_easy & 0.054 & 0.050 & 0.004 \\
MH\_02\_easy & 0.052 & 0.040 & 0.012 \\
MH\_03\_medium & 0.079 & 0.050 & 0.029 \\
\textcolor{red}{MH\_04\_difficult} & 0.152 & 0.120 & 0.032 \\
MH\_05\_difficult & 0.142 & 0.100 & 0.042 \\
V1\_01\_easy & 0.094 & 0.090 & 0.004 \\
\textcolor{red}{V1\_02\_medium} & 0.095 & 0.080 & 0.015 \\
\textcolor{red}{V1\_03\_difficult} & 0.670 & 0.170 & 0.500 \\
\textcolor{red}{V2\_01\_easy} & 0.086 & 0.100 & -0.014 \\
V2\_02\_medium & 0.204 & 0.100 & 0.104 \\
\bottomrule
\end{tabular}
\end{table}
\textbf{Note:} there is no official OV$^{2}$SLAM Fast Mode testing results, rather they tested the Fast mode on sped-up versions of the EuRoC sequences instead. Therefore, I used the OV$^{2}$SLAM no LC mode as the closest comparison to my Test OV$^{2}$SLAM Fast mode.

\begin{table}[H]
\centering
\small
\caption{Current ORB-SLAM3 keyframe RMSE measurements vs ORB-SLAM3 stereo paper rows.}
\label{tab:orb_vs_paper_accuracy}
\begin{tabular}{lrrr}
\toprule
Sequence & Measured ORB keyframe RMSE (m) & Paper ORB-SLAM3 stereo (m) & Delta (m) \\
\midrule
MH\_01\_easy & 0.043 & 0.029 & 0.014 \\
MH\_02\_easy & 0.071 & 0.019 & 0.052 \\
MH\_03\_medium & 0.247 & 0.024 & 0.223 \\
\textcolor{red}{MH\_04\_difficult} & 0.039 & 0.085 & -0.046 \\
MH\_05\_difficult & 0.091 & 0.052 & 0.039 \\
V1\_01\_easy & 0.131 & 0.035 & 0.096 \\
\textcolor{red}{V1\_02\_medium} & 0.331 & 0.025 & 0.306 \\
\textcolor{red}{V1\_03\_difficult} & 0.069 & 0.061 & 0.008 \\
\textcolor{red}{V2\_01\_easy} & n/a & 0.041 & n/a \\
V2\_02\_medium & 0.228 & 0.028 & 0.200 \\
\bottomrule
\end{tabular}
\end{table}

A significant problem that faced me when benchmarking ORB-SLAM3 was the instability of the runs. Some runs got very low RMSE (e.g. 0.01) which appears to be an amazing run, but in reality, ORB-SLAM only processed 4/84 seconds of a sequence. In other cases, the RMSE went as high as RMSE=1 and 0.9, and in some others it fluctuates between 0.24, 0.35, and 0.12. A possible explanation of this phenomenon are the 3 things below:
\begin{itemize}
    \item ORB-SLAM3 tracking being too expensive CPU-operations wise, for example it was not parallelized so it causes more delays and more frame dropping $\to$ higher RMSE (the problem with this is that I allowed ORB-SLAM3 to work on all cores) 
    \item ROS overhead is a possibility since the original code was never packaged with ROS and a community-maintained wrapper was downloaded to allow for ROS functionalities
\end{itemize}

\begin{table}[H]
\centering
\small
\caption{Current ORB-SLAM3 timing summary for the same sequences used in the paper-comparison table. Loop-closing means are taken from the end-of-run ORB log when present.}
\label{tab:orb_vs_paper_timing}
\begin{tabular}{lrrrr}
\toprule
Sequence & Track mean (ms) & Track P95 (ms) & LBA mean (ms) & Loop close mean (ms) \\
\midrule
MH\_01\_easy & 50.106 & 66.792 & 9.315 & n/a \\
MH\_02\_easy & 48.757 & 69.400 & 8.003 & n/a \\
MH\_03\_medium & 51.713 & 77.086 & 4.928 & 92.345 \\
\textcolor{red}{MH\_04\_difficult} & 38.323 & 59.090 & 4.115 & 175.540 \\
MH\_05\_difficult & 46.074 & 80.017 & 5.702 & 108.900 \\
V1\_01\_easy & 32.325 & 37.799 & 2.746 & n/a \\
\textcolor{red}{V1\_02\_medium} & 37.758 & 45.515 & 1.437 & n/a \\
\textcolor{red}{V1\_03\_difficult} & 46.427 & 65.468 & 1.696 & n/a \\
\textcolor{red}{V2\_01\_easy} & n/a & n/a & n/a & n/a \\
V2\_02\_medium & 36.495 & 45.114 & 3.871 & 49.983 \\
\bottomrule
\end{tabular}
\end{table}

\begin{figure}[H]
\centering
\includegraphics[width=\textwidth]{images/orbslam_cpu_by_sequence.png}
\caption{ORB-SLAM3 mean process CPU utilization by sequence with standard deviation error bars and value labels.}
\label{fig:orb_cpu_by_sequence}
\end{figure}

From the figure above, we can observe that ORB-SLAM3 can reach up to a maximum of \~217+\% CPU utilization, which suggests that it should be run on 3 cores to avoid contention but current experiements show that with all 4 cores, the algorithm is not fully stable.

\clearpage

\section{Logical Thread-role Tables}

\begin{table}[H]
\centering
\scriptsize
\caption{OV$^{2}$SLAM logical thread-role mean utilization aggregated across sequences.}
\label{tab:ov2_role_means_table}
\begin{tabular}{lrrr}
\toprule
Role & Accurate mean (\%) & Fast mean (\%) & Delta (fast-accurate) \\
\midrule
ov2\_est & 63.386 & 14.445 & -48.941 \\
ov2\_lc & 42.399 & n/a & n/a \\
ov2\_main & 14.343 & 6.192 & -8.150 \\
ov2\_map & 12.944 & 4.509 & -8.435 \\
ov2\_feed & 1.085 & 0.626 & -0.459 \\
OV$^{2}$SLAM\_node\_main & 1.118 & 0.801 & -0.318 \\
ov2\_merge & 0.000 & 0.000 & 0.000 \\
ov2\_vizfr & 0.000 & 0.000 & 0.000 \\
ov2\_vizkf & 0.000 & 0.000 & 0.000 \\
\bottomrule
\end{tabular}
\end{table}

\begin{table}[H]
\centering
\scriptsize
\caption{OV$^{2}$SLAM logical thread-role utilization variability across sequences.}
\label{tab:ov2_role_std_table}
\begin{tabular}{lrr}
\toprule
Role & Accurate std across sequences (\%) & Fast std across sequences (\%) \\
\midrule
ov2\_est & 6.659 & 2.293 \\
ov2\_lc & 11.259 & n/a \\
ov2\_main & 0.982 & 0.786 \\
ov2\_map & 3.184 & 1.274 \\
ov2\_feed & 0.073 & 0.021 \\
OV$^{2}$SLAM\_node\_main & 0.080 & 0.041 \\
ov2\_merge & 0.000 & 0.000 \\
ov2\_vizfr & 0.000 & 0.000 \\
ov2\_vizkf & 0.000 & 0.000 \\
\bottomrule
\end{tabular}
\end{table}

\begin{table}[H]
\centering
\small
\caption{ORB-SLAM3 logical thread-role utilization aggregated across completed sequences.}
\label{tab:orb_role_table}
\begin{tabular}{lrrr}
\toprule
Role & Mean CPU (\%) & Std across sequences (\%) & Sequences present \\
\midrule
orb\_gba & 78.657 & 8.950 & 4 \\
orb\_lmap & 20.356 & 14.882 & 9 \\
orb\_track & 19.150 & 1.803 & 9 \\
stereo\_node\_cpp & 2.892 & 0.225 & 9 \\
orb\_loop & 2.082 & 1.561 & 9 \\
recvUC & 1.068 & 0.079 & 9 \\
\bottomrule
\end{tabular}
\end{table}
As seen and as expected, bundle adjustment is the highest CPU-consuming thread.

\begin{figure}[H]
\centering
\includegraphics[width=\textwidth]{images/orbslam_role_means.png}
\caption{ORB-SLAM3 logical thread-role utilization aggregated across completed sequences.}
\label{fig:orb_roles}
\end{figure}

\clearpage

\section{OV$^{2}$SLAM Timing Delta Table}

\begin{longtable}{llrrrr}
\caption{OV$^{2}$SLAM accurate vs fast timing deltas by sequence and timer.}\label{tab:ov2_timing_deltas}\\
\toprule
Sequence & Timer & Accurate (ms) & Fast (ms) & Delta & \% change \\
\midrule
\endfirsthead
\toprule
Sequence & Timer & Accurate (ms) & Fast (ms) & Delta & \% change \\
\midrule
\endhead
\bottomrule
\endfoot
MH\_01\_easy & Full FE & 17.083 & 4.895 & $\downarrow$ 12.188 & -71.3\% \\
MH\_01\_easy & FE Track & 15.028 & 4.383 & $\downarrow$ 10.645 & -70.8\% \\
MH\_01\_easy & FE KF create & 16.872 & 3.763 & $\downarrow$ 13.109 & -77.7\% \\
MH\_01\_easy & Mapper KF proc & 47.601 & 8.310 & $\downarrow$ 39.291 & -82.5\% \\
MH\_01\_easy & Local BA & 292.162 & 41.283 & $\downarrow$ 250.879 & -85.9\% \\
MH\_01\_easy & LC ProcessKF & 82.758 & n/a & n/a & n/a \\
MH\_02\_easy & Full FE & 17.088 & 5.025 & $\downarrow$ 12.064 & -70.6\% \\
MH\_02\_easy & FE Track & 14.631 & 4.427 & $\downarrow$ 10.204 & -69.7\% \\
MH\_02\_easy & FE KF create & 17.054 & 3.876 & $\downarrow$ 13.178 & -77.3\% \\
MH\_02\_easy & Mapper KF proc & 51.340 & 8.460 & $\downarrow$ 42.881 & -83.5\% \\
MH\_02\_easy & Local BA & 280.729 & 38.392 & $\downarrow$ 242.338 & -86.3\% \\
MH\_02\_easy & LC ProcessKF & 83.636 & n/a & n/a & n/a \\
MH\_03\_medium & Full FE & 19.522 & 5.088 & $\downarrow$ 14.434 & -73.9\% \\
MH\_03\_medium & FE Track & 15.628 & 4.313 & $\downarrow$ 11.315 & -72.4\% \\
MH\_03\_medium & FE KF create & 20.002 & 3.800 & $\downarrow$ 16.202 & -81.0\% \\
MH\_03\_medium & Mapper KF proc & 57.868 & 8.497 & $\downarrow$ 49.371 & -85.3\% \\
MH\_03\_medium & Local BA & 211.707 & 33.928 & $\downarrow$ 177.779 & -84.0\% \\
MH\_03\_medium & LC ProcessKF & 64.739 & n/a & n/a & n/a \\
\textcolor{red}{MH\_04\_difficult} & Full FE & 15.578 & 4.990 & $\downarrow$ 10.589 & -68.0\% \\
\textcolor{red}{MH\_04\_difficult} & FE Track & 12.727 & 4.170 & $\downarrow$ 8.557 & -67.2\% \\
\textcolor{red}{MH\_04\_difficult} & FE KF create & 16.121 & 4.220 & $\downarrow$ 11.902 & -73.8\% \\
\textcolor{red}{MH\_04\_difficult} & Mapper KF proc & 31.519 & 6.961 & $\downarrow$ 24.558 & -77.9\% \\
\textcolor{red}{MH\_04\_difficult} & Local BA & 195.941 & 28.558 & $\downarrow$ 167.383 & -85.4\% \\
\textcolor{red}{MH\_04\_difficult} & LC ProcessKF & 64.356 & n/a & n/a & n/a \\
MH\_05\_difficult & Full FE & 15.815 & 5.036 & $\downarrow$ 10.779 & -68.2\% \\
MH\_05\_difficult & FE Track & 13.087 & 4.273 & $\downarrow$ 8.814 & -67.3\% \\
MH\_05\_difficult & FE KF create & 16.426 & 4.229 & $\downarrow$ 12.197 & -74.3\% \\
MH\_05\_difficult & Mapper KF proc & 33.355 & 6.370 & $\downarrow$ 26.985 & -80.9\% \\
MH\_05\_difficult & Local BA & 223.581 & 31.544 & $\downarrow$ 192.037 & -85.9\% \\
MH\_05\_difficult & LC ProcessKF & 66.218 & n/a & n/a & n/a \\
V1\_01\_easy & Full FE & 19.733 & 5.114 & $\downarrow$ 14.618 & -74.1\% \\
V1\_01\_easy & FE Track & 16.799 & 4.309 & $\downarrow$ 12.490 & -74.4\% \\
V1\_01\_easy & FE KF create & 20.179 & 4.086 & $\downarrow$ 16.093 & -79.8\% \\
V1\_01\_easy & Mapper KF proc & 71.542 & 8.346 & $\downarrow$ 63.195 & -88.3\% \\
V1\_01\_easy & Local BA & 271.416 & 32.296 & $\downarrow$ 239.120 & -88.1\% \\
V1\_01\_easy & LC ProcessKF & 75.109 & n/a & n/a & n/a \\
\textcolor{red}{V1\_02\_medium} & Full FE & 21.698 & 5.698 & $\downarrow$ 16.001 & -73.7\% \\
\textcolor{red}{V1\_02\_medium} & FE Track & 14.990 & 4.156 & $\downarrow$ 10.834 & -72.3\% \\
\textcolor{red}{V1\_02\_medium} & FE KF create & 20.409 & 4.315 & $\downarrow$ 16.094 & -78.9\% \\
\textcolor{red}{V1\_02\_medium} & Mapper KF proc & 43.554 & 7.464 & $\downarrow$ 36.090 & -82.9\% \\
\textcolor{red}{V1\_02\_medium} & Local BA & 140.398 & 17.691 & $\downarrow$ 122.707 & -87.4\% \\
\textcolor{red}{V1\_02\_medium} & LC ProcessKF & 73.415 & n/a & n/a & n/a \\
\textcolor{red}{V1\_03\_difficult} & Full FE & 21.225 & 6.105 & $\downarrow$ 15.120 & -71.2\% \\
\textcolor{red}{V1\_03\_difficult} & FE Track & 13.773 & 3.886 & $\downarrow$ 9.887 & -71.8\% \\
\textcolor{red}{V1\_03\_difficult} & FE KF create & 20.282 & 5.953 & $\downarrow$ 14.329 & -70.6\% \\
\textcolor{red}{V1\_03\_difficult} & Mapper KF proc & 35.529 & 4.934 & $\downarrow$ 30.595 & -86.1\% \\
\textcolor{red}{V1\_03\_difficult} & Local BA & 86.274 & 7.846 & $\downarrow$ 78.428 & -90.9\% \\
\textcolor{red}{V1\_03\_difficult} & LC ProcessKF & 66.747 & n/a & n/a & n/a \\
\textcolor{red}{V2\_01\_easy} & Full FE & 15.033 & 4.908 & $\downarrow$ 10.125 & -67.4\% \\
\textcolor{red}{V2\_01\_easy} & FE Track & 12.450 & 4.093 & $\downarrow$ 8.357 & -67.1\% \\
\textcolor{red}{V2\_01\_easy} & FE KF create & 16.287 & 4.733 & $\downarrow$ 11.554 & -70.9\% \\
\textcolor{red}{V2\_01\_easy} & Mapper KF proc & 34.761 & 7.099 & $\downarrow$ 27.662 & -79.6\% \\
\textcolor{red}{V2\_01\_easy} & Local BA & 212.102 & 28.772 & $\downarrow$ 183.330 & -86.4\% \\
\textcolor{red}{V2\_01\_easy} & LC ProcessKF & 68.419 & n/a & n/a & n/a \\
V2\_02\_medium & Full FE & 20.488 & 5.793 & $\downarrow$ 14.695 & -71.7\% \\
V2\_02\_medium & FE Track & 14.194 & 4.038 & $\downarrow$ 10.156 & -71.6\% \\
V2\_02\_medium & FE KF create & 18.857 & 5.329 & $\downarrow$ 13.529 & -71.7\% \\
V2\_02\_medium & Mapper KF proc & 42.090 & 10.646 & $\downarrow$ 31.444 & -74.7\% \\
V2\_02\_medium & Local BA & 107.253 & 16.232 & $\downarrow$ 91.021 & -84.9\% \\
V2\_02\_medium & LC ProcessKF & 74.974 & n/a & n/a & n/a \\
\end{longtable}
The task with the highest average delay is Local BA (local bundle adjustment). In the accurate stereo configuration, it has the largest mean runtime of around 202.16 ms, which is much higher than the front-end tracking stages by about 187ms. This is normal because the local bundle adjustment is an optimization algorithm that refines the camera poses and map points, and it is much more computationally expensive. Interestingly, the same operation is being used in the Front-end thread according to the paper, however, this Front-end thread is only a small motion-only PnP refinement of the current pose with fixed map points and a very tight number of iterations, while the actual local bundle adjustment optimizes multiple keyframes and map points which takes significantly more time.
\clearpage

% ═══════════════════════════════════════════════════════════════════════════════
\section{RTAB-Map Results}
% ═══════════════════════════════════════════════════════════════════════════════

\subsection{RTAB-Map RMSE Summary}

The RTAB-Map stereo odometry + SLAM pipeline was evaluated across all ten EuRoC sequences. Each sequence was repeated three times; results are reported as \textbf{mean$\pm$std across all valid runs}, where a run is valid if it has RMSE $>0.01$\,m and $\geq$50 recorded poses (low-RMSE, low-pose runs correspond to early tracking loss producing trivially small errors). The evaluation metric is Absolute Pose Error (APE) with Sim(3) Umeyama alignment using \texttt{evo\_ape}.

\begin{table}[H]
\centering
\small
\caption{RTAB-Map mean$\pm$std RMSE per sequence (across valid runs) vs.\ OV$^2$SLAM Accurate and ORB-SLAM3 (both best run). n/a = algorithm did not complete the sequence.}
\label{tab:rtabmap_rmse}
\begin{tabular}{lrrrr}
\toprule
Sequence & RTAB-Map avg (m) & Valid runs & OV2 acc (m) & ORB-SLAM3 (m) \\
\midrule
MH\_01\_easy       & $0.102 \pm 0.063$ & 3/3 & 0.040 & 0.043 \\
MH\_02\_easy       & $0.190 \pm 0.143$ & 3/3 & 0.049 & 0.071 \\
MH\_03\_medium     & $0.146 \pm 0.038$ & 3/3 & 0.046 & 0.247 \\
\textcolor{red}{MH\_04\_difficult} & $0.224 \pm 0.070$ & 3/3 & 0.069 & 0.039 \\
MH\_05\_difficult  & $0.166$ (2/3 lost tracking) & 1/3 & 0.067 & 0.091 \\
V1\_01\_easy       & $0.096 \pm 0.029$ & 3/3 & 0.055 & 0.131 \\
\textcolor{red}{V1\_02\_medium}   & $0.081$ (2/3 lost tracking) & 1/3 & 0.088 & 0.331 \\
\textcolor{red}{V1\_03\_difficult}& $0.133 \pm 0.046$ & 3/3 & 0.118 & 0.069 \\
\textcolor{red}{V2\_01\_easy}     & $0.161 \pm 0.052$ & 3/3 & 0.072 & n/a \\
V2\_02\_medium     & $0.244 \pm 0.035$ & 3/3 & 0.054 & 0.228 \\
\bottomrule
\end{tabular}
\end{table}

\begin{figure}[H]
\centering
\includegraphics[width=\textwidth]{images/rtabmap_rmse_per_sequence.png}
\caption{RTAB-Map mean APE RMSE per sequence with $\pm$1\,std error bars across valid runs. Green bars = easy sequences; red = medium/difficult sequences. Sequences with only one valid run show no error bar. Dashed lines mark 5\,cm and 10\,cm reference thresholds.}
\label{fig:rtabmap_rmse}
\end{figure}

Comparing the three systems: RTAB-Map's average RMSE is consistently higher than OV$^2$SLAM Accurate. Against ORB-SLAM3 the picture is mixed: ORB-SLAM3 wins on sequences where it runs successfully (e.g.\ MH\_01: 0.043 vs.\ 0.102\,m; V1\_03: 0.069 vs.\ 0.133\,m; MH\_04: 0.039 vs.\ 0.224\,m), but completely fails on \textcolor{red}{V1\_02\_medium} (0.331\,m vs.\ RTAB-Map's 0.081\,m) and does not run on \textcolor{red}{V2\_01\_easy}. RTAB-Map is the only system that completes every sequence. The large standard deviations (e.g.\ MH\_02: $\pm$0.143\,m) reflect the hardware bottleneck: at 8.4--10.5\,Hz effective odometry, outcomes are sensitive to OS scheduling and which frames happen to be queued, making results non-deterministic across runs.

\subsection{RTAB-Map Stereo Odometry Timing}

The \texttt{stereo\_odometry} node runs the GFTT+FLANN feature-based tracking pipeline for each incoming stereo frame. Timing is measured from the \texttt{OdomInfo.time\_estimation} and \texttt{local\_bundle\_time} fields published on \texttt{/odom\_info}. Values are averaged across all valid runs per sequence.

\begin{table}[H]
\centering
\small
\caption{RTAB-Map stereo odometry estimation time per sequence (average over valid runs).}
\label{tab:rtabmap_timing_odom}
\begin{tabular}{lrrr}
\toprule
Sequence & Mean (ms) & Std (ms) & P95 (ms) \\
\midrule
MH\_01\_easy       & 101.7 & 29.9 & 160.0 \\
MH\_02\_easy       & 104.4 & 28.5 & 168.3 \\
MH\_03\_medium     & 97.2  & 28.3 & 137.7 \\
MH\_04\_difficult  & 104.2 & 30.3 & 157.6 \\
MH\_05\_difficult  & 88.3  & 24.8 & 133.3 \\
V1\_01\_easy       & 103.6 & 30.0 & 161.3 \\
V1\_02\_medium     & 93.0  & 27.6 & 147.8 \\
V1\_03\_difficult  & 92.8  & 27.5 & 135.0 \\
V2\_01\_easy       & 100.3 & 27.2 & 153.2 \\
V2\_02\_medium     & 90.4  & 23.4 & 135.5 \\
\midrule
\textbf{Mean}      & \textbf{97.6} & \textbf{27.7} & \textbf{149.0} \\
\bottomrule
\end{tabular}
\end{table}

\begin{table}[H]
\centering
\small
\caption{RTAB-Map local bundle adjustment time per sequence (average over valid runs). Triggered selectively on keyframe creation; zero when no keyframe is created.}
\label{tab:rtabmap_timing_lb}
\begin{tabular}{lrr}
\toprule
Sequence & Mean (ms) & Std (ms) \\
\midrule
MH\_01\_easy       & 25.6 & 12.8 \\
MH\_02\_easy       & 24.7 & 12.9 \\
MH\_03\_medium     & 13.7 & 14.8 \\
MH\_04\_difficult  & 25.8 & 17.2 \\
MH\_05\_difficult  & 22.3 & 13.8 \\
V1\_01\_easy       & 23.0 & 13.0 \\
V1\_02\_medium     & 15.5 & 16.5 \\
V1\_03\_difficult  & 9.3  & 15.1 \\
V2\_01\_easy       & 22.4 & 13.6 \\
V2\_02\_medium     & 8.2  & 11.2 \\
\midrule
\textbf{Mean}      & \textbf{19.1} & \textbf{14.1} \\
\bottomrule
\end{tabular}
\end{table}

\begin{figure}[H]
\centering
\includegraphics[width=\textwidth]{images/rtabmap_odom_timing_per_sequence.png}
\caption{RTAB-Map stereo odometry timing per sequence. Blue bars = mean odometry estimation time ($\pm$1\,std); orange bars = mean local bundle time ($\pm$1\,std); red triangles ($\blacktriangle$) = P95 estimation time. The red dashed line marks 50\,ms (20\,Hz budget); the orange dotted line marks 100\,ms (10\,Hz).}
\label{fig:rtabmap_timing}
\end{figure}

A critical observation is that the mean odometry estimation time \textbf{consistently exceeds 88--104\,ms across all sequences} (std $\approx$\,25--30\,ms, P95 = 133--162\,ms), which is \textbf{more than double the 50\,ms budget} imposed by the 20\,Hz stereo stream. This means that RTAB-Map's \texttt{stereo\_odometry} cannot process every incoming frame in real time on the Raspberry Pi 5. With \texttt{always\_process\_most\_recent\_frame = false}, the node queues incoming frames, gradually accumulating a backlog. The effective odometry processing rate was verified directly from the \texttt{/odom\_info} timestamps recorded during each run: the node processes at \textbf{8.4--10.5\,Hz} across all sequences, roughly half of the 20\,Hz input rate. This means approximately every other frame is eventually processed, with increasing temporal lag relative to the camera. This frame lag introduces growing misalignment between the motion model and the actual camera trajectory, which is the primary contributor to the elevated RMSE on the more dynamic and difficult sequences.

The P95 processing time (133--162\,ms) is especially important: it means that at peak load, a single odometry call takes over 3$\times$ the frame budget, causing the backlog to deepen temporarily. These spikes coincide with keyframe creation events where the local bundle adjustment (orange bars, mean 8--26\,ms, std 11--18\,ms) runs on top of the feature estimation, as seen from the wide standard deviation of the local bundle time.

\subsection{RTAB-Map CPU Utilization}

CPU usage was monitored using a custom \texttt{monitor\_cpu.py} sampler running at 1\,Hz throughout each sequence. CPU percentages are measured as $\Delta$(user+system time) / $\Delta$(wall time) per thread, where 100\% represents full utilization of one CPU core. The monitored pipeline consists of two ROS2 nodes: \texttt{stereo\_odometry} and \texttt{rtabmap}.

\begin{table}[H]
\centering
\small
\caption{RTAB-Map CPU utilization per sequence (average over valid runs, pipeline = stereo\_odometry + rtabmap).}
\label{tab:rtabmap_cpu}
\begin{tabular}{lrrr}
\toprule
Sequence & stereo\_odometry (\%) & rtabmap (\%) & Combined (\%) \\
\midrule
MH\_01\_easy       & 80.7  & 51.7 & 132.3 \\
MH\_02\_easy       & 86.6  & 48.0 & 134.6 \\
MH\_03\_medium     & 88.0  & 37.3 & 125.3 \\
MH\_04\_difficult  & 87.3  & 42.7 & 130.0 \\
MH\_05\_difficult  & 85.8  & 45.4 & 131.2 \\
V1\_01\_easy       & 83.4  & 47.7 & 131.1 \\
V1\_02\_medium     & 80.0  & 40.1 & 120.1 \\
V1\_03\_difficult  & 87.2  & 33.1 & 120.3 \\
V2\_01\_easy       & 87.0  & 41.2 & 128.2 \\
V2\_02\_medium     & 89.5  & 33.5 & 123.0 \\
\midrule
\textbf{Mean}      & \textbf{85.6} & \textbf{42.1} & \textbf{127.6} \\
\bottomrule
\end{tabular}
\end{table}

\begin{figure}[H]
\centering
\includegraphics[width=\textwidth]{images/rtabmap_cpu_per_sequence.png}
\caption{RTAB-Map average CPU utilization per sequence, broken down by node. Stacked bars show \texttt{stereo\_odometry} (blue) and \texttt{rtabmap} (red); black error bars indicate $\pm$1\,std of the combined total across all 1\,Hz CPU samples. Labels show mean$\pm$std. Dashed lines mark 100\% (one core) and 200\% (two cores).}
\label{fig:rtabmap_cpu_per_sequence}
\end{figure}

The combined RTAB-Map pipeline consumes an average of \textbf{127.6\%} CPU, meaning it consistently requires slightly more than one full CPU core across the entire Raspberry Pi 5's four-core budget. The \texttt{stereo\_odometry} node dominates at \textbf{85.6\%}, driven by the GFTT feature detection, FLANN matching, and RANSAC outlier rejection — all of which internally exploit OpenCV's parallel primitives across the available cores. The \texttt{rtabmap} node contributes a further \textbf{42.1\%} on average, reflecting the SLAM graph update, working memory management, and loop closure detection running on a separate thread.

\textbf{Note:} these measurements cover only the two RTAB-Map nodes. The additional pipeline components (two \texttt{rectify\_node} instances for stereo rectification, \texttt{euroc\_caminfo\_pub.py}, and \texttt{ros2 bag play}) are not included and would raise the total system CPU consumption further.

\begin{figure}[H]
\centering
\includegraphics[width=0.55\textwidth]{images/rtabmap_cpu_per_thread.png}
\caption{RTAB-Map average CPU per node, averaged across all sequences and valid runs ($\pm$1\,std across sequences shown as error bars). DDS transport threads excluded. Y-axis extends to 400\% for comparison with the four-core budget.}
\label{fig:rtabmap_cpu_per_thread}
\end{figure}

\subsection{RTAB-Map Paper Comparison}

The reference results are taken from Table~7 of Labb{\'e} \& Michaud (2024) (\textit{RTAB-Map as an Open-Source Lidar and Visual SLAM Library}, arXiv:2403.06341), which reports ATE in centimetres for each odometry strategy on EuRoC stereo. The row used is \textbf{F2M} (Frame-to-Map), which corresponds to the \texttt{OdomF2M} strategy configured in this experiment. The paper results were obtained on a desktop-class CPU running the odometry node only (no SLAM/loop-closure node). By contrast, the measurements here include the full pipeline: F2M stereo odometry \emph{plus} the \texttt{rtabmap} SLAM node with appearance-based loop closure. Loop closure can improve accuracy on revisited areas but adds CPU cost that further loads the Raspberry Pi 5.

\begin{table}[H]
\centering
\small
\caption{RTAB-Map mean RMSE (averaged across valid runs) vs.\ F2M odometry-only results from Labb{\'e} \& Michaud (2024), Table~7 (desktop CPU, no loop closure).}
\label{tab:rtabmap_vs_paper}
\begin{tabular}{lrrr}
\toprule
Sequence & Measured avg RMSE (m) & Paper F2M ATE (m) & Delta avg (m) \\
\midrule
MH\_01\_easy       & $0.102 \pm 0.063$ & 0.017 & +0.085 \\
MH\_02\_easy       & $0.190 \pm 0.143$ & 0.025 & +0.165 \\
MH\_03\_medium     & $0.146 \pm 0.038$ & 0.068 & +0.078 \\
MH\_04\_difficult  & $0.224 \pm 0.070$ & 0.160 & +0.064 \\
MH\_05\_difficult  & $0.166$ (1/3 runs valid) & 0.076 & +0.090 \\
V1\_01\_easy       & $0.096 \pm 0.029$ & 0.071 & +0.025 \\
V1\_02\_medium     & $0.081$ (1/3 runs valid) & 0.040 & +0.041 \\
V1\_03\_difficult  & $0.133 \pm 0.046$ & 0.097 & +0.036 \\
V2\_01\_easy       & $0.161 \pm 0.052$ & 0.082 & +0.079 \\
V2\_02\_medium     & $0.244 \pm 0.035$ & 0.120 & +0.124 \\
\bottomrule
\end{tabular}
\end{table}

Several observations stand out:
\begin{itemize}
    \item The gap is now larger than it appeared with best-run values, because the average captures the full variability across runs. The high std on MH\_02 ($\pm$0.143\,m) reflects one catastrophically drifting run (RMSE = 0.390\,m) alongside two good ones.
    \item The tightest sequences relative to the paper are \textbf{V1\_01} (+25\,mm), \textbf{V1\_02} (+41\,mm), and \textbf{V1\_03} (+36\,mm), suggesting loop closure is most effective in the Vicon Room environment.
    \item The largest gaps are on \textbf{MH\_02} (+0.165\,m) and \textbf{V2\_02\_medium} (+0.124\,m), where both frame lag and scene complexity cause run-to-run inconsistency.
    \item The comparison is inherently unfair: the paper runs odometry-only at 20\,Hz on a desktop, while the measured results run full SLAM at 8.4--10.5\,Hz on a 4-core ARM.
\end{itemize}

\clearpage

\section{Per-sequence Trajectories}

This section collects the XY trajectory overlays for OV2SLAM accurate, OV2SLAM fast, and ORB-SLAM3 for each available EuRoC sequence.

\subsection{MH\_01\_easy}

\begin{figure}[H]
\centering
\begin{minipage}[t]{0.24\textwidth}
\centering
\includegraphics[width=\linewidth]{images/ov2_accurate_MH_01_easy_trajectory_xy.png}
\\
\footnotesize OV2 accurate
\end{minipage}
\hfill
\begin{minipage}[t]{0.24\textwidth}
\centering
\includegraphics[width=\linewidth]{images/ov2_fast_MH_01_easy_trajectory_xy.png}
\\
\footnotesize OV2 fast
\end{minipage}
\hfill
\begin{minipage}[t]{0.24\textwidth}
\centering
\includegraphics[width=\linewidth]{images/orbslam3_MH_01_easy_trajectory_xy.png}
\\
\footnotesize ORB-SLAM3
\end{minipage}
\hfill
\begin{minipage}[t]{0.24\textwidth}
\centering
\includegraphics[width=\linewidth]{images/rtabmap_MH_01_easy_trajectory_xy.png}
\\
\footnotesize RTAB-Map (run 3)
\end{minipage}
\caption{MH\_01\_easy XY trajectory overlays: OV2SLAM accurate, OV2SLAM fast, ORB-SLAM3, and RTAB-Map best run.}
\label{fig:trajectory_MH_01_easy}
\end{figure}

\subsection{MH\_02\_easy}

\begin{figure}[H]
\centering
\begin{minipage}[t]{0.24\textwidth}
\centering
\includegraphics[width=\linewidth]{images/ov2_accurate_MH_02_easy_trajectory_xy.png}
\\
\footnotesize OV2 accurate
\end{minipage}
\hfill
\begin{minipage}[t]{0.24\textwidth}
\centering
\includegraphics[width=\linewidth]{images/ov2_fast_MH_02_easy_trajectory_xy.png}
\\
\footnotesize OV2 fast
\end{minipage}
\hfill
\begin{minipage}[t]{0.24\textwidth}
\centering
\includegraphics[width=\linewidth]{images/orbslam3_MH_02_easy_trajectory_xy.png}
\\
\footnotesize ORB-SLAM3
\end{minipage}
\hfill
\begin{minipage}[t]{0.24\textwidth}
\centering
\includegraphics[width=\linewidth]{images/rtabmap_MH_02_easy_trajectory_xy.png}
\\
\footnotesize RTAB-Map (run 2)
\end{minipage}
\caption{MH\_02\_easy XY trajectory overlays: OV2SLAM accurate, OV2SLAM fast, ORB-SLAM3, and RTAB-Map best run.}
\label{fig:trajectory_MH_02_easy}
\end{figure}

\subsection{MH\_03\_medium}

\begin{figure}[H]
\centering
\begin{minipage}[t]{0.24\textwidth}
\centering
\includegraphics[width=\linewidth]{images/ov2_accurate_MH_03_medium_trajectory_xy.png}
\\
\footnotesize OV2 accurate
\end{minipage}
\hfill
\begin{minipage}[t]{0.24\textwidth}
\centering
\includegraphics[width=\linewidth]{images/ov2_fast_MH_03_medium_trajectory_xy.png}
\\
\footnotesize OV2 fast
\end{minipage}
\hfill
\begin{minipage}[t]{0.24\textwidth}
\centering
\includegraphics[width=\linewidth]{images/orbslam3_MH_03_medium_trajectory_xy.png}
\\
\footnotesize ORB-SLAM3
\end{minipage}
\hfill
\begin{minipage}[t]{0.24\textwidth}
\centering
\includegraphics[width=\linewidth]{images/rtabmap_MH_03_medium_trajectory_xy.png}
\\
\footnotesize RTAB-Map (run 1)
\end{minipage}
\caption{MH\_03\_medium XY trajectory overlays: OV2SLAM accurate, OV2SLAM fast, ORB-SLAM3, and RTAB-Map best run.}
\label{fig:trajectory_MH_03_medium}
\end{figure}

\subsection{\textcolor{red}{MH\_04\_difficult}}

\begin{figure}[H]
\centering
\begin{minipage}[t]{0.24\textwidth}
\centering
\includegraphics[width=\linewidth]{images/ov2_accurate_MH_04_difficult_trajectory_xy.png}
\\
\footnotesize OV2 accurate
\end{minipage}
\hfill
\begin{minipage}[t]{0.24\textwidth}
\centering
\includegraphics[width=\linewidth]{images/ov2_fast_MH_04_difficult_trajectory_xy.png}
\\
\footnotesize OV2 fast
\end{minipage}
\hfill
\begin{minipage}[t]{0.24\textwidth}
\centering
\includegraphics[width=\linewidth]{images/orbslam3_MH_04_difficult_trajectory_xy.png}
\\
\footnotesize ORB-SLAM3
\end{minipage}
\hfill
\begin{minipage}[t]{0.24\textwidth}
\centering
\includegraphics[width=\linewidth]{images/rtabmap_MH_04_difficult_trajectory_xy.png}
\\
\footnotesize RTAB-Map (run 3)
\end{minipage}
\caption{\textcolor{red}{MH\_04\_difficult} XY trajectory overlays: OV2SLAM accurate, OV2SLAM fast, ORB-SLAM3, and RTAB-Map best run.}
\label{fig:trajectory_MH_04_difficult}
\end{figure}

\subsection{MH\_05\_difficult}

\begin{figure}[H]
\centering
\begin{minipage}[t]{0.24\textwidth}
\centering
\includegraphics[width=\linewidth]{images/ov2_accurate_MH_05_difficult_trajectory_xy.png}
\\
\footnotesize OV2 accurate
\end{minipage}
\hfill
\begin{minipage}[t]{0.24\textwidth}
\centering
\includegraphics[width=\linewidth]{images/ov2_fast_MH_05_difficult_trajectory_xy.png}
\\
\footnotesize OV2 fast
\end{minipage}
\hfill
\begin{minipage}[t]{0.24\textwidth}
\centering
\includegraphics[width=\linewidth]{images/orbslam3_MH_05_difficult_trajectory_xy.png}
\\
\footnotesize ORB-SLAM3
\end{minipage}
\hfill
\begin{minipage}[t]{0.24\textwidth}
\centering
\includegraphics[width=\linewidth]{images/rtabmap_MH_05_difficult_trajectory_xy.png}
\\
\footnotesize RTAB-Map (run 1)
\end{minipage}
\caption{MH\_05\_difficult XY trajectory overlays: OV2SLAM accurate, OV2SLAM fast, ORB-SLAM3, and RTAB-Map best run.}
\label{fig:trajectory_MH_05_difficult}
\end{figure}

\subsection{V1\_01\_easy}

\begin{figure}[H]
\centering
\begin{minipage}[t]{0.24\textwidth}
\centering
\includegraphics[width=\linewidth]{images/ov2_accurate_V1_01_easy_trajectory_xy.png}
\\
\footnotesize OV2 accurate
\end{minipage}
\hfill
\begin{minipage}[t]{0.24\textwidth}
\centering
\includegraphics[width=\linewidth]{images/ov2_fast_V1_01_easy_trajectory_xy.png}
\\
\footnotesize OV2 fast
\end{minipage}
\hfill
\begin{minipage}[t]{0.24\textwidth}
\centering
\includegraphics[width=\linewidth]{images/orbslam3_V1_01_easy_trajectory_xy.png}
\\
\footnotesize ORB-SLAM3
\end{minipage}
\hfill
\begin{minipage}[t]{0.24\textwidth}
\centering
\includegraphics[width=\linewidth]{images/rtabmap_V1_01_easy_trajectory_xy.png}
\\
\footnotesize RTAB-Map (run 3)
\end{minipage}
\caption{V1\_01\_easy XY trajectory overlays: OV2SLAM accurate, OV2SLAM fast, ORB-SLAM3, and RTAB-Map best run.}
\label{fig:trajectory_V1_01_easy}
\end{figure}

\subsection{\textcolor{red}{V1\_02\_medium}}

\begin{figure}[H]
\centering
\begin{minipage}[t]{0.24\textwidth}
\centering
\includegraphics[width=\linewidth]{images/ov2_accurate_V1_02_medium_trajectory_xy.png}
\\
\footnotesize OV2 accurate
\end{minipage}
\hfill
\begin{minipage}[t]{0.24\textwidth}
\centering
\includegraphics[width=\linewidth]{images/ov2_fast_V1_02_medium_trajectory_xy.png}
\\
\footnotesize OV2 fast
\end{minipage}
\hfill
\begin{minipage}[t]{0.24\textwidth}
\centering
\includegraphics[width=\linewidth]{images/orbslam3_V1_02_medium_trajectory_xy.png}
\\
\footnotesize ORB-SLAM3
\end{minipage}
\hfill
\begin{minipage}[t]{0.24\textwidth}
\centering
\includegraphics[width=\linewidth]{images/rtabmap_V1_02_medium_trajectory_xy.png}
\\
\footnotesize RTAB-Map (run 3)
\end{minipage}
\caption{\textcolor{red}{V1\_02\_medium} XY trajectory overlays: OV2SLAM accurate, OV2SLAM fast, ORB-SLAM3, and RTAB-Map best run.}
\label{fig:trajectory_V1_02_medium}
\end{figure}

\subsection{\textcolor{red}{V1\_03\_difficult}}

\begin{figure}[H]
\centering
\begin{minipage}[t]{0.24\textwidth}
\centering
\includegraphics[width=\linewidth]{images/ov2_accurate_V1_03_difficult_trajectory_xy.png}
\\
\footnotesize OV2 accurate
\end{minipage}
\hfill
\begin{minipage}[t]{0.24\textwidth}
\centering
\includegraphics[width=\linewidth]{images/ov2_fast_V1_03_difficult_trajectory_xy.png}
\\
\footnotesize OV2 fast
\end{minipage}
\hfill
\begin{minipage}[t]{0.24\textwidth}
\centering
\includegraphics[width=\linewidth]{images/orbslam3_V1_03_difficult_trajectory_xy.png}
\\
\footnotesize ORB-SLAM3
\end{minipage}
\hfill
\begin{minipage}[t]{0.24\textwidth}
\centering
\includegraphics[width=\linewidth]{images/rtabmap_V1_03_difficult_trajectory_xy.png}
\\
\footnotesize RTAB-Map (run 2)
\end{minipage}
\caption{\textcolor{red}{V1\_03\_difficult} XY trajectory overlays: OV2SLAM accurate, OV2SLAM fast, ORB-SLAM3, and RTAB-Map best run.}
\label{fig:trajectory_V1_03_difficult}
\end{figure}

\subsection{\textcolor{red}{V2\_01\_easy}}

\begin{figure}[H]
\centering
\begin{minipage}[t]{0.24\textwidth}
\centering
\includegraphics[width=\linewidth]{images/ov2_accurate_V2_01_easy_trajectory_xy.png}
\\
\footnotesize OV2 accurate
\end{minipage}
\hfill
\begin{minipage}[t]{0.24\textwidth}
\centering
\includegraphics[width=\linewidth]{images/ov2_fast_V2_01_easy_trajectory_xy.png}
\\
\footnotesize OV2 fast
\end{minipage}
\hfill
\begin{minipage}[t]{0.24\textwidth}
\centering
\fbox{\footnotesize ORB-SLAM3 n/a}
\end{minipage}
\hfill
\begin{minipage}[t]{0.24\textwidth}
\centering
\includegraphics[width=\linewidth]{images/rtabmap_V2_01_easy_trajectory_xy.png}
\\
\footnotesize RTAB-Map (run 2)
\end{minipage}
\caption{\textcolor{red}{V2\_01\_easy} XY trajectory overlays: OV2SLAM accurate, OV2SLAM fast, ORB-SLAM3 (not available), and RTAB-Map best run.}
\label{fig:trajectory_V2_01_easy}
\end{figure}

\subsection{V2\_02\_medium}

\begin{figure}[H]
\centering
\begin{minipage}[t]{0.24\textwidth}
\centering
\includegraphics[width=\linewidth]{images/ov2_accurate_V2_02_medium_trajectory_xy.png}
\\
\footnotesize OV2 accurate
\end{minipage}
\hfill
\begin{minipage}[t]{0.24\textwidth}
\centering
\includegraphics[width=\linewidth]{images/ov2_fast_V2_02_medium_trajectory_xy.png}
\\
\footnotesize OV2 fast
\end{minipage}
\hfill
\begin{minipage}[t]{0.24\textwidth}
\centering
\includegraphics[width=\linewidth]{images/orbslam3_V2_02_medium_trajectory_xy.png}
\\
\footnotesize ORB-SLAM3
\end{minipage}
\hfill
\begin{minipage}[t]{0.24\textwidth}
\centering
\includegraphics[width=\linewidth]{images/rtabmap_V2_02_medium_trajectory_xy.png}
\\
\footnotesize RTAB-Map (run 3)
\end{minipage}
\caption{V2\_02\_medium XY trajectory overlays: OV2SLAM accurate, OV2SLAM fast, ORB-SLAM3, and RTAB-Map best run.}
\label{fig:trajectory_V2_02_medium}
\end{figure}

\clearpage


\section{Conclusions}

While OV$^{2}$SLAM was benchmarked to the best of my ability, ORB-SLAM3 presents a significant challenge, particularly because it is not available on the official ROS GitHub and is unstable during runtime. My hypothesis is that
\begin{itemize}
    \item \textbf{Hypothesis 1: Pipeline Latency.} The rosbag adapter publishes stereo frames at real-time rate (\texttt{playback\_rate: 1.0}), but ROS2 topic transport and subscriber callback overhead MAY have introduced sufficient latency that by the time each stereo pair reaches the ORB-SLAM3 tracking thread, the subsequent frame has already been published. The bounded stereo buffer (maximum 20 frames per channel) then makes non-deterministic keep/drop decisions under load, causing different frames to survive across runs and producing run-to-run RMSE variability.

    \item \textbf{Hypothesis 2: Tracking Budget Exceeded.} The ORB-SLAM3 tracking thread exceeds the 50\,ms per-frame budget imposed by EuRoC's 20\,Hz stereo input. When per-frame tracking time surpasses this threshold, the pipeline accumulates a backlog faster than it can be consumed. The bounded buffer subsequently drops frames non-deterministically depending on OS thread scheduling, causing divergent map evolution paths and unstable trajectory estimates across repeated runs.
\end{itemize}

\subsection*{RTAB-Map Summary and Four-Way Comparison}

RTAB-Map presents a well-integrated ROS2 pipeline and completed every run without crashing — a key advantage over ORB-SLAM3. The system is limited primarily by hardware: the stereo odometry pipeline takes 88--104\,ms per frame (P95 up to 162\,ms), more than double the 50\,ms budget at 20\,Hz. The effective processing rate was verified from \texttt{/odom\_info} timestamps at \textbf{8.4--10.5\,Hz} — roughly every other frame is processed with growing temporal lag.

\paragraph{RMSE comparison (averaged across valid runs).}
\begin{itemize}
    \item \textbf{OV$^2$SLAM Accurate} achieves the lowest average RMSE (mean $\approx$\,0.066\,m), tracking all sequences reliably with results close to its published paper values.
    \item \textbf{OV$^2$SLAM Fast} is worse than Accurate on harder sequences (mean $\approx$\,0.106\,m excl.\ V1\_03) and catastrophically fails on \textcolor{red}{V1\_03\_difficult} (0.670\,m).
    \item \textbf{ORB-SLAM3} delivers competitive best-run results on sequences it completes (0.039--0.247\,m), but shows severe instability: it fails entirely on \textcolor{red}{V2\_01\_easy} (n/a), produces 0.331\,m on \textcolor{red}{V1\_02\_medium}, and has high run-to-run variability (average across successful sequences $\approx$\,0.139\,m). Unlike RTAB-Map, its failures are not always detectable from pose count alone.
    \item \textbf{RTAB-Map} has the highest mean RMSE ($\approx$\,0.154\,m) but is the only system completing all sequences. It beats ORB-SLAM3 on \textcolor{red}{V1\_02\_medium} (0.081 vs.\ 0.331\,m) and MH\_03\_medium (0.146 vs.\ 0.247\,m). Its failure mode — early tracking loss producing a low pose count — is easy to detect and flag automatically.
\end{itemize}

\paragraph{Frontend/odometry latency comparison.}
\begin{itemize}
    \item \textbf{OV$^2$SLAM Fast}: $\approx$\,5\,ms mean — the fastest, comfortably within the 50\,ms budget.
    \item \textbf{OV$^2$SLAM Accurate}: $\approx$\,16--22\,ms mean — still well within the 50\,ms budget.
    \item \textbf{ORB-SLAM3}: $\approx$\,32--52\,ms mean — at or slightly above the budget, causing non-deterministic frame drops.
    \item \textbf{RTAB-Map}: \textbf{88--104\,ms mean} ($\pm$25--30\,ms std, P95 up to 162\,ms) — consistently over 2$\times$ the budget, forcing operation at 8.4--10.5\,Hz.
\end{itemize}

\paragraph{CPU utilization comparison.}
\begin{itemize}
    \item \textbf{OV$^2$SLAM Fast}: $\approx$\,43--52\% mean — the lightest, easily fitting in a single core.
    \item \textbf{RTAB-Map}: $\approx$\,120--135\% mean — uses just over one full core; the most CPU-efficient \emph{complete SLAM system} evaluated.
    \item \textbf{ORB-SLAM3}: $\approx$\,87--150\% mean — similar to RTAB-Map but less stable.
    \item \textbf{OV$^2$SLAM Accurate}: $\approx$\,141--165\% mean — highest average CPU of the four, driven by its local bundle adjustment and loop closure threads.
\end{itemize}

\paragraph{Overall assessment.} On a Raspberry Pi 5, OV$^2$SLAM Fast is the clear winner for edge deployment: lowest latency, lowest CPU, and acceptable accuracy. OV$^2$SLAM Accurate is the best choice when accuracy matters and CPU budget allows $\approx$1.6 cores. RTAB-Map's odometry is too slow for 20\,Hz real-time operation on this hardware but its failure mode is graceful and its loop closure can partially compensate for drift. ORB-SLAM3 remains the least suitable due to non-deterministic instability.

%% ============================================================
\section{Pairwise Comparison}
\label{sec:pairwise}
%% ============================================================

This section presents side-by-side bar-chart comparisons across all four algorithms — ORB-SLAM3, OV$^2$SLAM~Accurate, OV$^2$SLAM~Fast, and RTAB-Map — evaluated on the same ten EuRoC stereo sequences.
Each plot shows the mean value per sequence, with error bars representing $\pm1$ standard deviation.
For RTAB-Map the statistics are computed across valid runs (RMSE $>1\,$cm and $\geq50$ matched poses); for OV$^2$SLAM and ORB-SLAM3 each sequence was run once (no std available, bars without error caps).

% -------------------------------------------------------
\subsection{Accuracy (APE RMSE)}
\label{sec:pairwise_rmse}
% -------------------------------------------------------

\begin{figure}[htbp]
  \centering
  \includegraphics[width=\linewidth]{images/pairwise_rmse.png}
  \caption{%
    \textbf{Pairwise APE RMSE per sequence.}
    OV$^2$SLAM~Accurate is the most consistently accurate algorithm, staying below 12\,cm on all sequences.
    OV$^2$SLAM~Fast degrades noticeably on harder sequences (V1\_03, MH\_04, V2\_02).
    RTAB-Map (mean $\pm$ std across valid runs) achieves 8--24\,cm and is competitive with OV$^2$SLAM~Accurate on V1\_01 and V1\_02.
    ORB-SLAM3 shows the highest variability; it failed entirely on V2\_01.
    Dashed horizontal lines mark 5\,cm and 10\,cm thresholds.
  }
  \label{fig:pairwise_rmse}
\end{figure}

\begin{table}[htbp]
  \centering
  \caption{Pairwise APE RMSE summary (m). RTAB-Map reports mean $\pm$ std; others single run. ``---'' = algorithm did not complete.}
  \label{tab:pairwise_rmse}
  \small
  \begin{tabular}{lcccc}
    \hline
    Sequence & ORB-SLAM3 & OV$^2$SLAM Acc. & OV$^2$SLAM Fast & RTAB-Map \\
    \hline
    MH\_01\_easy      & 0.043 & 0.040 & 0.054 & $0.102 \pm 0.000$ \\
    MH\_02\_easy      & 0.071 & 0.049 & 0.052 & $0.190 \pm 0.059$ \\
    MH\_03\_medium    & 0.247 & 0.046 & 0.079 & $0.146 \pm 0.036$ \\
    MH\_04\_difficult & 0.039 & 0.069 & 0.152 & $0.223 \pm 0.000$ \\
    MH\_05\_difficult & 0.091 & 0.067 & 0.142 & $0.166 \pm 0.000$ \\
    V1\_01\_easy      & 0.131 & 0.055 & 0.094 & $0.096 \pm 0.036$ \\
    V1\_02\_medium    & 0.331 & 0.088 & 0.095 & $0.081 \pm 0.000$ \\
    V1\_03\_difficult & 0.069 & 0.118 & 0.670 & $0.133 \pm 0.000$ \\
    V2\_01\_easy      & ---   & 0.072 & 0.086 & $0.161 \pm 0.000$ \\
    V2\_02\_medium    & 0.228 & 0.054 & 0.204 & $0.244 \pm 0.093$ \\
    \hline
  \end{tabular}
\end{table}

% -------------------------------------------------------
\subsection{CPU Utilisation}
\label{sec:pairwise_cpu}
% -------------------------------------------------------

\begin{figure}[htbp]
  \centering
  \includegraphics[width=\linewidth]{images/pairwise_cpu.png}
  \caption{%
    \textbf{Pairwise CPU utilisation per sequence.}
    Values represent mean process CPU usage as a percentage of one core (the device has 4 cores, so 400\% = full saturation).
    OV$^2$SLAM~Fast is by far the lightest ($\approx$43--52\%), using less than half a core.
    RTAB-Map ($\approx$120--135\%) and ORB-SLAM3 ($\approx$87--150\%) are comparable.
    OV$^2$SLAM~Accurate is the heaviest ($\approx$142--165\%), driven by its local bundle adjustment and loop closure threads.
    Error bars show $\pm1$ std of the per-sample CPU distribution within each sequence.
    Dashed lines mark 100\,\% (one core) and 200\,\% (two cores).
  }
  \label{fig:pairwise_cpu}
\end{figure}

\begin{table}[htbp]
  \centering
  \caption{Pairwise process CPU utilisation (mean $\pm$ std, \% of one core). V2\_01 ORB-SLAM3 did not complete (---).}
  \label{tab:pairwise_cpu}
  \small
  \begin{tabular}{lcccc}
    \hline
    Sequence & ORB-SLAM3 & OV$^2$SLAM Acc. & OV$^2$SLAM Fast & RTAB-Map \\
    \hline
    MH\_01\_easy      & $140 \pm 61$ & $153 \pm 87$  & $43 \pm 35$  & $132 \pm 18$ \\
    MH\_02\_easy      & $135 \pm 62$ & $150 \pm 81$  & $44 \pm 36$  & $135 \pm 17$ \\
    MH\_03\_medium    & $146 \pm 66$ & $161 \pm 90$  & $48 \pm 38$  & $125 \pm 17$ \\
    MH\_04\_difficult & $150 \pm 80$ & $147 \pm 88$  & $43 \pm 34$  & $130 \pm 18$ \\
    MH\_05\_difficult & $144 \pm 70$ & $147 \pm 87$  & $42 \pm 34$  & $131 \pm 18$ \\
    V1\_01\_easy      & $87  \pm 51$ & $157 \pm 88$  & $46 \pm 36$  & $131 \pm 15$ \\
    V1\_02\_medium    & $112 \pm 52$ & $165 \pm 104$ & $47 \pm 36$  & $120 \pm 17$ \\
    V1\_03\_difficult & $132 \pm 61$ & $162 \pm 98$  & $47 \pm 32$  & $120 \pm 17$ \\
    V2\_01\_easy      & ---          & $142 \pm 83$  & $42 \pm 33$  & $128 \pm 16$ \\
    V2\_02\_medium    & $119 \pm 53$ & $164 \pm 95$  & $52 \pm 37$  & $123 \pm 17$ \\
    \hline
  \end{tabular}
\end{table}

% -------------------------------------------------------
\subsection{Frontend / Odometry Latency}
\label{sec:pairwise_latency}
% -------------------------------------------------------

\begin{figure}[htbp]
  \centering
  \includegraphics[width=\linewidth]{images/pairwise_latency.png}
  \caption{%
    \textbf{Pairwise per-frame processing latency per sequence.}
    Metric definitions differ by algorithm to reflect the most relevant per-frame cost:
    OV$^2$SLAM = full front-end timer;
    ORB-SLAM3 = ROS wrapper tracking call;
    RTAB-Map = stereo odometry estimation time.
    Error bars show $\pm1$ std.
    Dashed lines mark 50\,ms (20\,Hz real-time budget) and 100\,ms (10\,Hz).
    OV$^2$SLAM~Fast ($\approx$5\,ms) is the fastest by an order of magnitude.
    OV$^2$SLAM~Accurate ($\approx$15--22\,ms) remains comfortably within budget.
    ORB-SLAM3 ($\approx$32--52\,ms) is at or slightly over the 50\,ms line.
    RTAB-Map ($\approx$88--104\,ms) consistently exceeds 50\,ms, confirming sub-real-time operation on this hardware.
    V2\_01 ORB-SLAM3 failed and is shown as~0.
  }
  \label{fig:pairwise_latency}
\end{figure}

\begin{table}[htbp]
  \centering
  \caption{Pairwise per-frame processing latency (mean $\pm$ std, ms). See caption of Fig.~\ref{fig:pairwise_latency} for metric definitions.}
  \label{tab:pairwise_latency}
  \small
  \begin{tabular}{lcccc}
    \hline
    Sequence & ORB-SLAM3 & OV$^2$SLAM Acc. & OV$^2$SLAM Fast & RTAB-Map \\
    \hline
    MH\_01\_easy      & $50.1 \pm 9.2$   & $17.1 \pm 8.1$  & $4.9 \pm 1.5$  & $73.4 \pm 25.4$ \\
    MH\_02\_easy      & $48.8 \pm 9.9$   & $17.1 \pm 8.8$  & $5.0 \pm 1.5$  & $104.0 \pm 27.1$ \\
    MH\_03\_medium    & $51.7 \pm 13.7$  & $19.5 \pm 12.3$ & $5.1 \pm 1.7$  & $96.7 \pm 26.8$ \\
    MH\_04\_difficult & $38.3 \pm 11.5$  & $15.6 \pm 7.9$  & $5.0 \pm 1.9$  & $104.1 \pm 27.2$ \\
    MH\_05\_difficult & $46.1 \pm 18.8$  & $15.8 \pm 8.6$  & $5.0 \pm 1.9$  & $88.3 \pm 24.1$ \\
    V1\_01\_easy      & $32.3 \pm 3.3$   & $19.7 \pm 11.5$ & $5.1 \pm 1.8$  & $103.6 \pm 28.1$ \\
    V1\_02\_medium    & $37.8 \pm 4.8$   & $21.7 \pm 12.2$ & $5.7 \pm 2.2$  & $93.0 \pm 25.9$ \\
    V1\_03\_difficult & $46.4 \pm 11.1$  & $21.2 \pm 12.9$ & $6.1 \pm 3.2$  & $92.8 \pm 25.5$ \\
    V2\_01\_easy      & ---              & $15.0 \pm 8.1$  & $4.9 \pm 1.9$  & $100.2 \pm 27.4$ \\
    V2\_02\_medium    & $36.5 \pm 5.1$   & $20.5 \pm 11.8$ & $5.8 \pm 2.6$  & $90.4 \pm 25.2$ \\
    \hline
  \end{tabular}
\end{table}

\subsection*{Pairwise Comparison Summary}

The three bar charts together paint a consistent picture of the accuracy--efficiency frontier on the Raspberry Pi~5:

\begin{itemize}
  \item \textbf{OV$^2$SLAM~Fast} dominates the efficiency axis: 5\,ms latency, $<$50\% CPU, with accuracy degrading only on the harder sequences.
        It is the only algorithm that leaves meaningful headroom on all four cores simultaneously.

  \item \textbf{OV$^2$SLAM~Accurate} provides the best accuracy across all ten sequences at moderate cost ($\approx$1.6 cores, 15--22\,ms).
        It is the preferred choice whenever a second core can be dedicated to the SLAM process.

  \item \textbf{ORB-SLAM3} delivers competitive accuracy on most sequences but shows two failure modes: catastrophic tracking loss on difficult indoor scenes (V1\_02, V1\_03, MH\_03) and a complete failure on V2\_01.
        Its 32--52\,ms latency is borderline for 20\,Hz and it cannot be relied upon without robust recovery logic.

  \item \textbf{RTAB-Map} is the heaviest in latency (88--104\,ms) but has the most \emph{predictable} CPU footprint ($\approx$1.3 cores, low std).
        Its SLAM accuracy (8--24\,cm) is acceptable for navigation-class applications, and its failure mode is soft (tracking degrades rather than crashing).
        It is unsuitable for 20\,Hz real-time odometry on this hardware without hardware acceleration.
\end{itemize}

\paragraph{Why ORB-SLAM3 fails more than RTAB-Map despite being faster per frame.}
A counter-intuitive result from the pairwise plots is that ORB-SLAM3, despite having a \emph{lower} per-frame tracking latency (32--52\,ms vs.\ 88--104\,ms for RTAB-Map), suffers far worse \emph{tracking reliability} — catastrophic losses on three sequences and a complete failure on V2\_01.
The explanation lies in architectural differences, not raw speed:

\begin{enumerate}
  \item \textbf{Descriptor matching vs.\ optical flow.}
        ORB-SLAM3 tracks every frame by matching ORB binary descriptors against a \emph{global sparse map}.
        Descriptor matching fails abruptly when the map becomes stale (too few keyframes nearby) or when image conditions change (motion blur, poor lighting).
        RTAB-Map's stereo odometry instead uses Lucas--Kanade optical flow from the previous frame, which degrades \emph{gradually} — it produces noisier poses before it loses track rather than a hard failure.

  \item \textbf{Aggressive keyframe culling under resource pressure.}
        ORB-SLAM3 prunes redundant keyframes continuously to bound memory.
        On the Raspberry Pi~5, the local mapping thread runs slower than on a desktop, so keyframe culling can strip the local map too aggressively before the next keyframe is inserted.
        The result is a sparse local map against which incoming frames cannot find enough inliers, triggering tracking loss.

  \item \textbf{Relocalization cost and cascading failure.}
        When ORB-SLAM3 loses tracking it enters a relocalization mode that queries a visual bag-of-words (DBoW2) database — an expensive operation.
        On this hardware the relocalization often does not complete before several more frames arrive, causing the input queue to back up and making recovery increasingly unlikely once triggered.

  \item \textbf{Global Bundle Adjustment (GBA) map jumps.}
        ORB-SLAM3 can trigger a full GBA after a successful loop closure.
        A GBA correction updates \emph{all} keyframe and map-point positions at once; if the correction is large it can temporarily invalidate the motion model and cause the tracker to report a sudden, large pose error that it cannot recover from on the next frame.

  \item \textbf{RTAB-Map's memory management acts as a safety valve.}
        RTAB-Map's working-memory mechanism (controlled by \texttt{Rtabmap/MemoryThr}) limits how many recent signatures are kept in RAM.
        Old signatures are transferred to long-term memory (the database) rather than being discarded.
        This means the system \emph{degrades in loop-closure ability} as memory fills up, but the odometry front-end continues uninterrupted — there is no catastrophic failure path analogous to ORB-SLAM3's tracking-loss state.
\end{enumerate}

In short: ORB-SLAM3's faster per-frame cost comes at the price of a brittle all-or-nothing tracking model, whereas RTAB-Map's higher latency is accompanied by a more resilient, flow-based odometry that fails softly.
On a resource-constrained embedded platform, this difference in \emph{failure mode} can matter more than the difference in mean latency.

\end{document}